\documentclass[11pt,a4paper]{article}
\usepackage[utf8]{inputenc}
\usepackage[T1]{fontenc}
\usepackage[french]{babel}
\usepackage{lmodern}
\usepackage[a4paper,margin=2cm]{geometry}
\usepackage[table]{xcolor}
\usepackage{textcomp}
\usepackage{graphicx}
\usepackage{array}
\usepackage{booktabs}
\usepackage{longtable}
\usepackage{enumitem}
\usepackage{listings}
\usepackage{tikz}
\usetikzlibrary{arrows.meta, positioning, fit, backgrounds}
\usepackage{titlesec}
\usepackage[hidelinks]{hyperref}

\definecolor{titleblue}{HTML}{15314B}
\definecolor{accent}{HTML}{1D6FB8}
\definecolor{codebg}{HTML}{F4F6FB}
\definecolor{codecmt}{HTML}{55607A}
\definecolor{codekw}{HTML}{1D6FB8}
\definecolor{codestr}{HTML}{1A7F4B}

\titleformat{\section}{\Large\bfseries\color{titleblue}}{\thesection}{0.6em}{}
\titleformat{\subsection}{\large\bfseries\color{accent}}{\thesubsection}{0.6em}{}

\lstset{
  basicstyle=\ttfamily\footnotesize,
  backgroundcolor=\color{codebg},
  commentstyle=\color{codecmt}\itshape,
  keywordstyle=\color{codekw}\bfseries,
  stringstyle=\color{codestr},
  breaklines=true, breakatwhitespace=false,
  frame=single, framerule=0pt, framesep=6pt,
  xleftmargin=6pt, xrightmargin=6pt,
  showstringspaces=false, columns=fullflexible, keepspaces=true,
  literate={é}{{\'e}}1 {è}{{\`e}}1 {à}{{\`a}}1 {â}{{\^a}}1 {ê}{{\^e}}1
           {î}{{\^i}}1 {ô}{{\^o}}1 {û}{{\^u}}1 {ç}{{\c c}}1 {É}{{\'E}}1,
}

% Neutralise la ponctuation "active" de babel-french (:, ;, !, ?) : on saisit
% nous-mêmes les espaces. Indispensable pour que TikZ (label={above:...}) et
% listings compilent sans conflit de catcode.
\frenchsetup{AutoSpacePunctuation=false}

\begin{document}

\begin{center}
  {\LARGE\bfseries\color{titleblue} Construire et manipuler un environnement IA avec RAG}\\[4pt]
  {\large Retrieval-Augmented Generation --- base vectorielle locale + LLM local}\\[6pt]
  {\small PostgreSQL/pgvector \textperiodcentered{} Ollama (gemma2:2b, nomic-embed-text)
   \textperiodcentered{} Docker \textperiodcentered{} Python 3.14}
\end{center}
\vspace{4pt}\hrule\vspace{10pt}

\section{Objectif}
Le but du projet est de construire, de bout en bout, un environnement d'intelligence
artificielle exploitant le principe du \textbf{RAG (Retrieval-Augmented Generation)} :
enrichir les reponses d'un grand modele de langage (LLM) avec une base de connaissance
personnalisee, afin qu'il reponde sur des donnees qu'il n'a jamais vues lors de son
entrainement. La demonstration s'appuie sur une entreprise \textbf{fictive}, HelioniX
SAS (onduleurs photovoltaiques), dont les documents internes servent de base de
connaissance.

\section{Environnement construit}
L'environnement repose sur la \textbf{virtualisation applicative (conteneurs Docker)}
plutot que sur une machine virtuelle complete : chaque outil tourne dans un conteneur
isole, orchestre par \texttt{docker compose}. Ce choix reduit l'empreinte et rend
l'environnement reproductible en une commande.

\begin{center}
\renewcommand{\arraystretch}{1.3}
\begin{tabular}{>{\raggedright\arraybackslash}p{2.6cm} >{\raggedright\arraybackslash}p{3.6cm} >{\raggedright\arraybackslash}p{8.4cm}}
\toprule
\textbf{Couche} & \textbf{Element} & \textbf{Detail} \\
\midrule
Machine hote & Windows 11 Pro (x86-64) & Execute Docker Desktop (backend WSL 2) et l'interpreteur Python \\
Virtualisation & Docker Desktop & Conteneurisation, orchestration via \texttt{docker-compose.yml} \\
Base vectorielle & conteneur \texttt{rag-db} & image \texttt{pgvector/pgvector:pg16} --- PostgreSQL 16 + extension \texttt{vector}, port \texttt{localhost:5433} \\
Moteur LLM & conteneur \texttt{rag-ollama} & image \texttt{ollama/ollama} --- embedding + generation, port \texttt{localhost:11434} \\
Modele d'embedding & \texttt{nomic-embed-text} & Texte vers vecteur de \textbf{768} dimensions \\
Modele de generation & \texttt{gemma2:2b} (Google) & LLM leger (environ 1,6 Go) adapte a une execution locale (GPU/CPU) \\
Scripts & Python 3.14 (venv) & \texttt{psycopg2-binary} (PostgreSQL) + \texttt{requests} (API Ollama) \\
\bottomrule
\end{tabular}
\end{center}

\subsection{Schema d'architecture}

\begin{center}
\resizebox{0.98\textwidth}{!}{%
\begin{tikzpicture}[
  font=\small,
  cont/.style={draw=blue!55, fill=blue!4, rounded corners, align=left, inner sep=5pt},
  py/.style={draw=orange!70, fill=orange!12, rounded corners, align=left, inner sep=5pt},
  ar/.style={-{Latex[length=2.2mm]}, thick, draw=black!65},
]
  \node[cont] (ollama) {\textbf{rag-ollama} \texttt{ollama/ollama} (:11434)\\[2pt]
        \footnotesize gemma2:2b \textperiodcentered{} nomic-embed-text (768d) \textperiodcentered{} GPU};
  \node[cont, below=11mm of ollama] (db) {\textbf{rag-db} \texttt{pgvector/pgvector:pg16} (:5433)\\[2pt]
        \footnotesize table documents(embedding VECTOR(768)) \textperiodcentered{} index HNSW};
  \begin{scope}[on background layer]
    \node[draw=blue!40, fill=blue!2, rounded corners, fit=(ollama)(db), inner sep=6mm,
          label={[blue!55]above:\textbf{Docker Desktop (WSL 2)}}] (docker) {};
  \end{scope}
  \node[py, left=26mm of docker] (py) {\textbf{Scripts Python}\\[2pt]
        \footnotesize venv 3.14\\ \footnotesize ingest.py \textperiodcentered{} query.py};
  \begin{scope}[on background layer]
    \node[draw=black!55, rounded corners, fit=(py)(docker), inner sep=8mm,
          label={above:\textbf{Machine hote --- Windows 11}}] (host) {};
  \end{scope}
  \draw[ar] (py.east|-ollama.west) -- node[above,font=\scriptsize]{REST :11434} (ollama.west);
  \draw[ar] (py.east|-db.west) -- node[above,font=\scriptsize]{SQL :5433} (db.west);
\end{tikzpicture}%
}
\end{center}

\begin{center}\footnotesize\itshape Figure 1 --- Architecture : scripts Python sur l'hote,
conteneurs Ollama et PostgreSQL/pgvector orchestres par Docker. Les scripts dialoguent
avec Ollama en REST (port 11434) et avec PostgreSQL en SQL (port 5433).\end{center}

\section{Workflow des traitements}
Le pipeline se decompose en deux phases : une phase d'\textbf{ingestion} (hors ligne,
executee une fois pour peupler la base) et une phase d'\textbf{interrogation} (rejouee a
chaque question de l'utilisateur).

\begin{center}
\resizebox{0.98\textwidth}{!}{%
\begin{tikzpicture}[
  font=\small, node distance=7mm,
  n/.style={draw=blue!55, fill=blue!4, rounded corners, align=center,
            minimum height=11mm, inner sep=3pt, text width=26mm},
  g/.style={n, draw=orange!70, fill=orange!12},
  ar/.style={-{Latex[length=2.2mm]}, thick, draw=black!65},
  lane/.style={font=\bfseries\footnotesize, anchor=west},
]
  % Lane A : ingestion
  \node[n] (d)  {Documents\\ \texttt{data/*.md}};
  \node[n, right=of d]  (ch) {Decoupage\\ chunks 600/100};
  \node[g, right=of ch] (em) {Embedding\\ nomic-embed-text};
  \node[n, right=of em] (ins){INSERT\\ pgvector};
  \draw[ar] (d)--(ch); \draw[ar] (ch)--(em); \draw[ar] (em)--(ins);
  \node[lane] at ([yshift=8mm]d.north west) {A. Ingestion (ingest.py) --- hors ligne};
  % Lane B : interrogation
  \node[n, below=18mm of d] (q)  {Question\\ utilisateur};
  \node[g, right=of q]  (qe) {Embedding\\ de la question};
  \node[n, right=of qe] (kn) {Recherche KNN\\ pgvector TOP\_K=4};
  \node[n, right=of kn] (pr) {Prompt enrichi\\ contexte + question};
  \draw[ar] (q)--(qe); \draw[ar] (qe)--(kn); \draw[ar] (kn)--(pr);
  \node[g, below=11mm of kn] (gen) {Generation LLM\\ gemma2:2b};
  \node[n, right=of gen] (rep) {Reponse\\ contextualisee};
  \draw[ar] (pr.south) |- (gen.east);
  \draw[ar] (gen)--(rep);
  \node[lane] at ([yshift=8mm]q.north west) {B. Interrogation (query.py) --- a chaque question};
\end{tikzpicture}%
}
\end{center}

\begin{center}\footnotesize\itshape Figure 2 --- Workflow : (A) ingestion des documents,
(B) interrogation contextualisee du LLM par recherche de similarite vectorielle.\end{center}

\newpage
\section{Procedure d'installation}
L'ensemble de l'environnement se deploie depuis le dossier du projet.
\begin{lstlisting}[language=bash]
# 1) Demarrer la base vectorielle et le moteur LLM
docker compose up -d
docker compose ps                      # rag-db "healthy", rag-ollama "up"

# 2) Telecharger les modeles dans Ollama
docker exec rag-ollama ollama pull nomic-embed-text    # embeddings (768 dim.)
docker exec rag-ollama ollama pull gemma2:2b           # generation

# 3) Environnement Python
python -m venv .venv
.venv\Scripts\activate                 # Windows
pip install -r projet_cyber_rag/requirements.txt

# 4) Ingestion puis interrogation
python scripts/ingest.py --reset
python scripts/query.py "Quelle est la garantie de l'onduleur HX-Solar 5000 ?"
python scripts/query.py --no-rag "..." # comparaison : LLM seul, sans la base
python scripts/demo.py                 # batterie de demonstration complete
\end{lstlisting}
Au premier demarrage, le script \texttt{db/init.sql} est joue automatiquement : il active
l'extension \texttt{vector}, cree la table \texttt{documents} et son index HNSW.

\section{Scripts developpes}
\subsection{Ingestion : \texttt{ingest.py}}
Lit les documents, les decoupe en fragments (\textit{chunks}) d'environ 600 caracteres avec
100 de recouvrement, calcule l'embedding de chaque fragment via Ollama et l'insere dans
pgvector.
\begin{lstlisting}[language=Python]
def chunk_text(text, size, overlap):
    text = text.strip()
    if not text:
        return []
    chunks, start, n = [], 0, len(text)
    while start < n:
        end = min(start + size, n)
        if end < n:                       # coupe de preference sur un saut de ligne
            window = text.rfind("\n", start + size // 2, end)
            if window > start:
                end = window
        chunk = text[start:end].strip()
        if chunk:
            chunks.append(chunk)
        if end >= n:                      # tout le texte consomme -> fin
            break
        start = max(end - overlap, start + 1)   # progression stricte

def embed(text):                          # embedding via Ollama (nomic-embed-text)
    r = requests.post(f"{OLLAMA_URL}/api/embeddings",
                      json={"model": EMBED_MODEL, "prompt": text}, timeout=120)
    return r.json()["embedding"]
\end{lstlisting}

\subsection{Interrogation : \texttt{query.py}}
Vectorise la question, recupere les \texttt{TOP\_K=4} fragments les plus proches (distance
cosine via l'operateur \texttt{<=>} de pgvector), construit un prompt contraint au
contexte puis genere la reponse avec \texttt{gemma2:2b}.
\begin{lstlisting}[language=Python]
def retrieve(question, k=TOP_K):
    qvec = embed(question)                        # 1) vectorise la question
    cur.execute("""
        SELECT source, content, 1 - (embedding <=> %s::vector) AS score
        FROM documents
        ORDER BY embedding <=> %s::vector          -- distance cosine (pgvector)
        LIMIT %s;""", (qvec, qvec, k))            # 2) TOP_K chunks les plus proches
    return cur.fetchall()

def answer(question):
    hits = retrieve(question)
    context = "\n\n".join(f"[Source: {s}]\n{c}" for s, c, _ in hits)
    prompt = f"{SYSTEM_PROMPT}\n\n### CONTEXTE\n{context}\n\n### QUESTION\n{question}"
    return generate(prompt)                        # 3) generation gemma2:2b (Ollama)
\end{lstlisting}

\newpage
\section{Demonstration : couples prompt / reponse}
Le tableau ci-dessous met en evidence la \textbf{bonne prise en compte de la base de
connaissance} : les reponses \guillemotleft{} Avec RAG \guillemotright{} citent des faits
precis issus des documents
HelioniX. Pour 3 questions, la reponse \guillemotleft{} Sans RAG \guillemotright{} (LLM seul) est
fournie en regard : le modele est alors incapable de repondre correctement, ce qui
demontre l'apport du RAG. La derniere question porte volontairement sur un sujet
\textbf{hors base} : le systeme repond qu'il ne dispose pas de l'information, prouvant
qu'il ne brode pas.

\renewcommand{\arraystretch}{1.25}
\begin{longtable}{|>{\raggedright\arraybackslash}p{0.4cm}
                  |>{\raggedright\arraybackslash}p{3.6cm}
                  |>{\raggedright\arraybackslash}p{2.2cm}
                  |>{\raggedright\arraybackslash}p{7.6cm}|}
\hline
\rowcolor{titleblue!12}
\textbf{\#} & \textbf{Prompt} & \textbf{Chunks} & \textbf{Reponse} \\ \hline
\endfirsthead
\hline
\rowcolor{titleblue!12}
\textbf{\#} & \textbf{Prompt} & \textbf{Chunks} & \textbf{Reponse} \\ \hline
\endhead
1 & Quelle est la garantie constructeur de l'onduleur HX-Solar 5000 ? & 3 source(s)\newline score 0.812 & \textbf{Avec RAG :} 12 ans\newline\newline \textbf{Sans RAG :} La garantie constructeur de l'onduleur HX-Solar 5000 est de \textbf{10 ans} sur le matériel et \textbf{5 ans} sur les composants électroniques. \\ \hline
2 & Que signifie le code d'erreur E-204 sur le HX-Solar 5000 ? & 3 source(s)\newline score 0.787 & \textbf{Avec RAG :} Le code d'erreur E-204 signale une surchauffe du transformateur : il faut vérifier la ventilation et contacter le SAV HelioniX si le défaut persiste.\newline\newline \textbf{Sans RAG :} Le code d'erreur E-204 sur le HX-Solar 5000 indique un \textbf{défaut de communication avec la centrale solaire}. \\ \hline
3 & Combien de jours de télétravail par semaine sont autorisés chez HelioniX, et quel jour est obligatoirement travaillé au bureau ? & 3 source(s)\newline score 0.803 & \textbf{Avec RAG :} HelioniX autorise jusqu'à 3 jours de télétravail par semaine pour les postes éligibles. Le mardi est un jour de présence obligatoire au bureau pour toutes les équipes (réunion hebdomadaire « SunSync »). \\ \hline
4 & Quel est le rendement européen du HX-Solar 5000 et sur quel port TCP fonctionne le protocole SunLink v3 ? & 2 source(s)\newline score 0.786 & \textbf{Avec RAG :} Le rendement européen du HX-Solar 5000 est de 97,8 \%. Le protocole SunLink v3 fonctionne sur le port TCP 8770. \\ \hline
5 & Qui est la présidente de HelioniX et quel était le chiffre d'affaires 2024 ? & 3 source(s)\newline score 0.757 & \textbf{Avec RAG :} Amélie Rouanet est la présidente de HelioniX. Le chiffre d'affaires 2024 était de 38,6 millions d'euros. \\ \hline
6 & Quelle est la valeur faciale du ticket restaurant et la prise en charge employeur chez HelioniX ? & 3 source(s)\newline score 0.778 & \textbf{Avec RAG :} La valeur faciale du ticket restaurant est de 11 \texteuro{} et la prise en charge employeur est de 6,60 \texteuro{}. \\ \hline
7 & Quelle est la capitale de l'Australie ? & 3 source(s)\newline score 0.597 & \textbf{Avec RAG :} Je ne dispose pas de cette information dans ma base de connaissance.\newline\newline \textbf{Sans RAG :} Canberra  \\ \hline
\end{longtable}

\section{Conclusion}
L'environnement repond aux trois objectifs du projet : \textbf{(1)} construction d'une
plateforme IA locale conteneurisee (Ollama + PostgreSQL/pgvector), \textbf{(2)} ingestion
d'un jeu de documents dans une base de connaissance vectorielle, et \textbf{(3)}
interrogation d'un LLM enrichi par cette base. La comparaison \guillemotleft{} avec / sans RAG \guillemotright{}
confirme que la base de connaissance est effectivement mobilisee et qu'elle ameliore
nettement la pertinence des reponses sur des donnees metier privees.

\end{document}
